\chapter{Dyskusja i zakończenie}

\section{Realizacja celu pracy}
Główny cel pracy --- stworzenie konfigurowalnego środowiska AutoML dla finansowych
potoków Numerai --- został zrealizowany na poziomie architektury, implementacji
i badania empirycznego. AlphaPulse łączy pobieranie danych, walidowany schemat YAML,
routing cech, modele, HPO, ocenę czasową, zapis proweniencji i eksport funkcji
predykcyjnej. Benchmark pełnoskalowy potwierdził wykonalność kompletnego przepływu,
a seria v5.3 pokazała, że framework potrafi porównywać heterogeniczne konfiguracje
według jednego, domenowo poprawnego protokołu.

Wartość rozwiązania nie sprowadza się do pojedynczego estymatora. Kluczowym produktem
jest kontrolowane środowisko eksperymentalne: konfiguracja opisuje zamiar badacza,
warstwa wykonawcza realizuje ten sam podział i scoring, a baza prób wiąże wynik z kodem,
danymi i ziarnem. Takie rozdzielenie odpowiada założeniom współczesnego AutoML,
jednocześnie dodając wymagania specyficzne dla danych finansowych
\cite{automl-survey,optuna,ray}.

\section{Interpretacja wyników modelowych}
LightGBM i CatBoost osiągnęły najwyższe mediany CORR Sharpe na holdoucie. Jest to
spójne z pozycją gradient boostingu jako silnej klasy modeli dla danych tabelarycznych
\cite{lightgbm,catboost}. Nie oznacza jednak uniwersalnej przewagi tych algorytmów.
Przestrzeń była próbkowana adaptacyjnie, a modele fundamentowe działały z ostrzejszymi
limitami pamięciowymi. Poprawna interpretacja brzmi zatem: w przyjętym budżecie i dla
celu \texttt{target\_ender\_60} najlepsze konfiguracje pochodziły z dwóch rodzin
boostingowych.

Umiarkowana korelacja rankingów holdout--validation potwierdza znaczenie separacji
czasowej. W danych o niskim stosunku sygnału do szumu wybór najlepszego punktu na
jednym przedziale łatwo przejmuje jego lokalną specyfikę. Purging ogranicza przeciek
wynikający z nakładania się horyzontu celu, ale nie eliminuje niestacjonarności rynku
\cite{lopez}. Zewnętrzny przedział 88 er pełni dlatego inną rolę niż holdout HPO:
nie steruje samplerem i pozwala ocenić transfer wybranych receptur.

Najbardziej praktycznym wynikiem jest front Pareto. Próba 22 maksymalizuje CORR na
validation przy dodatnim MMC, natomiast próba 4 maksymalizuje MMC przy nieco słabszym
CORR. Są to dwa uzasadnione cele wdrożeniowe. Pierwszy profil preferuje stabilną
zgodność z celem, drugi większy wkład niezależny od Meta Modelu
\cite{numerai-corr,numerai-mmc}. Zachowanie obu kandydatów jest bardziej informacyjne
niż arbitralne połączenie metryk w jeden wskaźnik.

\section{Odpowiedzi na pytania eksperymentalne}
\begin{enumerate}
    \item \textbf{Skalowalność i reprodukowalność.} Pełnoskalowy potok przetworzył
    2,75 mln wierszy uczących i 3,90 mln wierszy walidacyjnych w 31,5 min. Każde
    uruchomienie zapisuje konfigurację, hash, metryki i artefakty, a końcowe HPO
    dodatkowo utrwala skróty danych oraz wersję kodu.
    \item \textbf{Rodziny modeli.} W badanym budżecie najwyższe mediany holdoutu
    uzyskały CatBoost i LightGBM. Wynik ma charakter opisowy, ponieważ HPO nie jest
    zrównoważoną ablacją.
    \item \textbf{Transfer i kompromis CORR--MMC.} Zgodność dwóch przedziałów jest
    dodatnia, ale umiarkowana. Próby 22 i 4 tworzą front Pareto na niezależnym
    validation i reprezentują odpowiednio profil CORR oraz MMC.
    \item \textbf{Koszt i kompletność procesu.} Analiza obejmuje 27 kompletnych
    prób z medianą 127,2 s. Baza TrialDB zachowała 84,4\% uruchomionej puli jako
    kompletne rekordy możliwe do bezpośredniego porównania.
\end{enumerate}

Odpowiedzi te odnoszą się bezpośrednio do pomiarów dostępnych w artefaktach. Nie
przypisują systemowi redukcji czasu pracy człowieka ani przewagi AutoResearch,
ponieważ takie twierdzenia wymagałyby osobnego eksperymentu porównawczego.

\section{Znaczenie inżynierskie}
Projekt pokazuje, że wymagania domenowe można zamknąć w kontraktach oprogramowania.
Kolumna ery nie jest zwykłą cechą, lecz steruje podziałem; horyzont celu wyznacza
separację; oficjalne metryki są częścią ewaluatora; a format submisji jest sprawdzany
w czystym podprocesie. Dzięki temu ryzyka metodologiczne są ograniczane przez
architekturę, a nie tylko przez dyscyplinę osoby prowadzącej eksperyment.

Istotna jest również warstwa artefaktów. 	exttt{protocol.json} dokumentuje wejście
i reguły badania, 	exttt{trials.db} przechowuje stan każdej próby, a W\&B umożliwia
porównanie krzywych oraz diagnostyk. Razem tworzą ślad audytowy potrzebny zarówno
w pracy naukowej, jak i w iteracyjnym utrzymaniu modelu. Z kolei eksport
\texttt{predict.pkl} zamyka lukę między eksperymentem i inferencją live.

\section{Ograniczenia}
Najważniejszym ograniczeniem jest budżet końcowej serii: jedno ziarno, 27 kompletnych
konfiguracji i dziesięcioprocentowa próba danych uczących. Nie pozwala to oszacować
wariancji między seedami ani przedstawić przedziałów ufności dla rankingu. Adaptacyjne
próbkowanie TPE utrudnia także przyczynową interpretację różnic między rodzinami.

Drugie ograniczenie dotyczy selekcji czasowej. Holdout uczestniczy w funkcji celu,
więc jego maksimum jest obciążone wielokrotnym porównaniem. Niezależny validation
ogranicza to ryzyko, lecz 88 er nadal odpowiada jednemu okresowi rynkowemu. Wyniki
nie są prognozą rezultatu live ani podstawą do automatycznego stakingu.

Benchmark pełnoskalowy używa innej wersji danych i celu niż końcowe HPO. Dostarcza
dowodu skalowalności, ale nie stanowi baseline jakości. W pracy celowo nie
ekstrapolowano wartości CORR lub MMC z próbki 10\% na pełny zbiór: zależność jakości
od liczby wierszy nie musi być liniowa.

\section{Rekomendacja wdrożeniowa i dalszy rozwój}
Pierwszym kandydatem do ponownego treningu na pełnych danych v5.3 jest próba 22,
ponieważ łączy najwyższy CORR na niezależnym validation z dodatnim MMC. Próba 4
powinna zostać zachowana jako alternatywa nastawiona na unikalność sygnału. Po
treningu pełnoskalowym oba modele należy oceniać w kolejnych rundach bez zmiany
konfiguracji; dopiero taki monitoring pozwoli oszacować stabilność live.

Dalszy rozwój powinien objąć powtórzenia dla wielu ziaren, kroczące okna obejmujące
więcej reżimów rynkowych oraz kontrolowane ablacje routingu cech, selekcji wariancji
i neutralizacji. AutoResearch jest gotowym rozszerzeniem warstwy automatyzacji, ale
jego wartość należy mierzyć przy identycznym budżecie i tej samej przestrzeni co TPE.
Naturalnym krokiem produkcyjnym jest także automatyczny harmonogram retreningu,
walidacji artefaktu i monitorowania dryfu metryk per-era.

\section{Wnioski końcowe}
AlphaPulse realizuje pełny, reprodukowalny cykl eksperymentu dla Numerai: od danych
i deklaratywnej konfiguracji, przez czasowo poprawną walidację i HPO, po diagnostykę
i eksport. Empirycznie system obsłużył trening pełnoskalowy oraz wyłonił na danych
v5.3 dwa niezdominowane profile modelu. Najbardziej zrównoważona próba 22 osiągnęła
CORR Sharpe 0,4247 i MMC Sharpe 0,2467 na niezależnym validation, a próba 4 podniosła
MMC Sharpe do 0,3063.

Najważniejszym wkładem pracy jest połączenie poprawnej metodyki finansowej z
inżynierią eksperymentów. Wynik nie jest pojedynczą liczbą, lecz systemem, który
umożliwia odtwarzanie decyzji, uczciwe porównywanie konfiguracji i przejście od
hipotezy do gotowego artefaktu predykcyjnego. W tym sensie AlphaPulse spełnia cel
konfigurowalnego środowiska AutoML dla finansowych potoków Numerai.
